
\documentclass[lettersize,journal]{IEEEtran}

% --- 数学、符号、字体宏包 ---
\usepackage{amsmath,amsfonts}
\usepackage{bm}
\usepackage{mathrsfs}
\usepackage{amssymb}

% --- 算法宏包（按需注释/启用）---
\usepackage{algorithm}
\usepackage{algpseudocode}
%\usepackage{algorithmic}
\usepackage[algo2e,ruled,vlined]{algorithm2e}
%\usepackage[ruled,linesnumbered]{algorithm2e}
\renewcommand{\algorithmicrequire}{\textbf{Input:}}
\renewcommand{\algorithmicensure}{\textbf{Output:}}

% --- 表格 / 图片 ---
\usepackage{array}
\usepackage[caption=false,font=normalsize,labelfont=sf,textfont=sf]{subfig}
\usepackage{graphicx}
\usepackage{epstopdf}
\usepackage{epsfig}
\usepackage{booktabs}

% --- 版面 / 其他 ---
\usepackage{textcomp}
\usepackage{stfloats}
\usepackage{url}
\usepackage{verbatim}
\usepackage{color}
\usepackage{setspace}
\usepackage{adjustbox}
\usepackage[numbers,sort&compress]{natbib}

% 数学公式允许跨页断行
\allowdisplaybreaks[4]
% 断词（可选）
\hyphenation{op-tical net-works semi-conduc-tor IEEE-Xplore}

% --- 定理 / 结论环境 ---
\newtheorem{theorem}{Theorem}
\newtheorem{corollary}{Corollary}
\newtheorem{definition}{Definition}
\newtheorem{lemma}{Lemma}
\newtheorem{proposition}{Proposition}
\newtheorem{assumption}{Assumption}
\newtheorem{remark}{Remark}
\newtheorem{example}{Example}
\newtheorem{proof}{Proof}
\newtheorem{case}{case}
\newtheorem{problem}{Problem}

\begin{document}

% ======================================================================
%  题目 / 作者 / 基金 / 联系方式（此处填本次论文信息）
% ======================================================================
\title{Placeholder: Paper Title\\*[1em]{\LARGE (rewrite this title)}}

\author{Author One~\IEEEmembership{...}, Author Two,
  Corresponding Author~\IEEEmembership{...}
\thanks{Financial support placeholder.}%
\thanks{Affiliation and email placeholder.}%
}

\maketitle

% ======================================================================
%  Abstract / Keywords
% ======================================================================
\begin{abstract}
Abstract placeholder: problem, method, stability conclusion, simulation.
\end{abstract}

\begin{IEEEkeywords}
Keywords placeholder: control systems, multi-agent, ...
\end{IEEEkeywords}

% ======================================================================
%  I. Introduction —— 由"文献调研结果"填充
%   段落规划：研究背景与领域发展 → 现状问题 → 本文方法 → 贡献列表
% ======================================================================
\section{Introduction}
\IEEEPARstart{O}{}verview placeholder.

Related work placeholder (cite via \cite{}).

Contributions:
\begin{enumerate}
\item Contribution 1 placeholder.
\item Contribution 2 placeholder.
\item Contribution 3 placeholder.
\end{enumerate}

The rest is organized as follows. Section II ..., Section III ..., Section IV ..., Section V ...

% 记法占位（Notations），若有需可在此列 $\cdot$ 定义。

% ======================================================================
%  II. Preliminaries and Problem Formulation
%      Graph Theory / 系统模型 / 攻击模型 / 假设与引理
% ======================================================================
\section{Preliminaries and Problem Formulation}

\subsection{Graph Theory}
Graph theory background placeholder.\begin{assumption}Assumption placeholder.\end{assumption}
\begin{lemma}Lemma placeholder.\end{lemma}

\subsection{Problem Formulation}
System model placeholder.
\begin{align}
% 状态方程占位
\dot{x}_{ik,1} &= x_{ik,2} \nonumber\\
\dot{x}_{ik,2} &= u_{ik}+f_{ik}+d_{ik}
\label{eq:system}
\end{align}
where the symbols are defined here.

\begin{assumption}Assumption placeholder.\end{assumption}
\begin{remark}Remark placeholder.\end{remark}

% ======================================================================
%  III. Main Result —— 由"理论草稿"填充（本类论文重心）
%      DRESO 设计 → 观测器设计 → 控制器设计 → 稳定性分析
% ======================================================================
\section{Main Result}

{\subsection{DRESO Design}}
Observer/error system placeholder.
\begin{align}
% 观测器方程占位
\label{eq:observer}
\end{align}
\begin{lemma}Lemma placeholder.\end{lemma}

{\subsection{PTFDRO Design}}
Distributed observer placeholder.
\begin{theorem}
Theorem statement placeholder.
\label{thm:1}
\end{theorem}
\textit{Proof:} Proof placeholder.

{\subsection{Prescribed-time Optimal Controller Design}}
Controller design placeholder.
\textit{Step 1:} placeholder.
\begin{align}
% 代价函数 / HJB 方程 / 控制器占位
\label{eq:controller}
\end{align}
\textit{Step 2:} placeholder.

{\subsection{Stability Analysis}}
Stability analysis placeholder.
\begin{theorem}
Stability theorem placeholder.
\label{thm:2}
\end{theorem}
\textit{Proof:} Using the Lyapunov function $V_i(t)$ Proof placeholder.

% ======================================================================
%  IV. Simulation —— 由"实验代码跑出的图片"填充
% ======================================================================
\section{Simulation}
Simulation setup placeholder (topology, parameters, baselines). Figures:

\begin{figure}[h]
	\centering
	% \includegraphics[height=9cm]{你的实验图.png}
	\caption{Placeholder: experimental figure 1.}
	\label{fig:1}
\end{figure}

% ======================================================================
%  V. Conclusion
% ======================================================================
\section{Conclusion}
Conclusion placeholder.

% ======================================================================
%  References：generated from "paper_search/paper_validate metadata" 生成 .bib
% ======================================================================
\bibliographystyle{IEEEtran}
\bibliography{IEEEabrv,main}

\end{document}